% ============================================================
% preamble.tex —— 高考数学真题集 · 公共导言区
% 引擎：XeLaTeX（文档类 ctexbook 在 main.tex 中声明）
% ============================================================

% ---------- 数学 ----------
\usepackage{amsmath}
\usepackage{amssymb}
\usepackage{mathtools}

% ---------- 绘图 ----------
\usepackage{tikz}
\usetikzlibrary{shapes.geometric, arrows.meta, calc, patterns, angles, quotes, positioning}
\usepackage{pgfplots}
\pgfplotsset{compat=1.18}

% ---------- 版面 ----------
\usepackage{geometry}
\geometry{a4paper, margin=2.2cm, headheight=15pt}
\usepackage{enumitem}
\setlist{itemsep=2pt, topsep=4pt, parsep=0pt}

% 中文与行内公式混排的断行弹性：
% emergencystretch 允许段落末路不通时追加伸展量；
% 降低二元/关系符罚分，允许行内长公式在 = 、+ 等处断行（否则整段外溢）
\emergencystretch=2em
\binoppenalty=700
\relpenalty=500
\usepackage{booktabs}      % 三线表
\usepackage[most]{tcolorbox}

% ---------- 选择题选项对齐 ----------
% tasks 宏包：\begin{gkopts}{列数} \task ... \end{gkopts}
% 自动等宽分栏，选项天然对齐；列数由转换脚本按选项长度决定（4/2/1）
\usepackage{tasks}
\settasks{label={}, label-width=0pt, item-indent=1.2em,
          column-sep=1em, after-item-skip=2pt, before-skip=2pt, after-skip=4pt}
\NewTasksEnvironment[style=itemize]{gkopts}[\task](4)

% ---------- 页眉页脚 ----------
\usepackage{fancyhdr}
\pagestyle{fancy}
\fancyhf{}
\fancyhead[L]{\small\leftmark}
\fancyhead[R]{\small 高考数学真题集 2017--2026}
\fancyfoot[C]{\small —— 第 \thepage\ 页 ——}
\renewcommand{\headrulewidth}{0.4pt}
% \chaptermark 只留章标题、去掉"第 x 章"前缀，页眉更紧凑
\renewcommand{\chaptermark}[1]{\markboth{#1}{}}

% ---------- 超链接（须最后加载）----------
\usepackage{hyperref}
\hypersetup{colorlinks=true, linkcolor=blue!55!black,
            bookmarksnumbered=true, unicode=true,
            pdftitle={高考数学真题集 2017-2026（全国 I 卷系）}}

% ============================================================
% 自定义命令与环境
% ============================================================

% ---------- 试卷标识（每套卷 \input 前由 \gkpaper 重设，用于 label 命名空间）----------
\newcommand{\gkpaperid}{none}

% ---------- 题号 ----------
% 试题部分的题号：加粗编号 + 自动 \label{q:<卷id>:<题号>}
\newcommand{\gkti}[1]{\par\medskip\noindent
  \textbf{#1．}\phantomsection\label{q:\gkpaperid:#1}\ignorespaces}

% 解析部分的题头：小节标题 + 回链到试题页
\newcommand{\gksolhead}[2]{%
  \subsubsection*{第 #1 题（#2）}%
  \phantomsection\label{s:\gkpaperid:#1}%
  {\small\color{gray}（题目见第~\pageref{q:\gkpaperid:#1}~页）}\par\nopagebreak}

% ---------- 填空题空位 ----------
\newcommand{\gkblank}{\underline{\hspace{2.2em}}}

% ---------- 提示框 ----------
% 卷种提示 / 解析归属声明（醒目，浅黄）
\newtcolorbox{gknotice}{colback=yellow!8, colframe=orange!60!black,
  boxrule=0.5pt, arc=1.5pt, left=6pt, right=6pt, top=4pt, bottom=4pt,
  before skip=8pt, after skip=8pt}

% 方法提示 / 注（浅蓝，用于解析后的点评）
\newtcolorbox{gktip}{colback=blue!4, colframe=blue!35!black,
  boxrule=0.4pt, arc=1.5pt, left=6pt, right=6pt, top=4pt, bottom=4pt,
  before skip=6pt, after skip=6pt, fontupper=\small}

% ---------- 答案行 ----------
% 解析中独立成段的"答案：X"用带底色的行呈现
\newcommand{\gkanswer}[1]{\par\smallskip\noindent
  \tcbox[colback=green!6, colframe=green!35!black, boxrule=0.4pt,
         arc=1.5pt, left=4pt, right=4pt, top=2pt, bottom=2pt,
         nobeforeafter]{\textbf{#1}}\par\smallskip}

% ---------- 结构标题（不进目录，避免目录被几百个条目淹没）----------
\newcommand{\gkpart}[1]{\subsection*{#1}\nopagebreak}      % 一、选择题
% （一）必考题 / 解析小问标题：块式加粗段，可自动换行（\paragraph 的 run-in 版式会溢出）
\newcommand{\gksubpart}[1]{\par\medskip\noindent\textbf{#1}\par\nopagebreak\smallskip}

% 罗马数字Ⅰ-Ⅻ（U+2160）与圈号①②③（U+2460）划入 CJK 字符类，
% 用中文字体渲染——西文字体 Latin Modern 缺这些字形，否则整段静默丢字
\xeCJKDeclareCharClass{CJK}{"2160 -> "217F, "2460 -> "24FF}

% ---------- 整卷装配 ----------
% \gkpaper{标识}{章标题}{文件基名}
%   papers/<基名>.tex   = 试题
%   answers/<基名>_sol.tex = 答案与解析
\newcommand{\gkpaper}[3]{%
  \renewcommand{\gkpaperid}{#1}%
  \chapter{#2}
  \section{试题}
  \input{papers/#3}
  \clearpage
  \section{参考答案与详细解析}
  \input{answers/#3_sol}
  \clearpage}
